% ===== Block B17: 原文 p.161–162 | 附录（超几何函数） | 译手 tr-ch3b =====

% [原文 p.161]
本附录中我们给出超几何函数的基本结果。对任意复数 $\alpha$ 与任意非负整数 $n$，定义
\[
(\alpha)_n=\alpha(\alpha+1)(\alpha+2)\cdots(\alpha+n-1),
\]
其中
\[
(\alpha)_0=1.
\]
利用伽马函数（gamma function）可将 $(\alpha)_n$ 写成
\[
(\alpha)_n=\frac{\Gamma(\alpha+n)}{\Gamma(\alpha)}.
\]
超几何级数 $F(\alpha,\beta,\gamma;z)$ 由下式定义：
\begin{equation}
F(\alpha,\beta,\gamma;z)=\sum_{n=0}^{\infty}\frac{(\alpha)_n(\beta)_n}{(\gamma)_nn!}z^n=1+\frac{\alpha\beta}{\gamma}z+\frac{\alpha(\alpha+1)\beta(\beta+1)}{\gamma(\gamma+1)2!}z^2+\cdots,\tag{A.39}
\end{equation}
其中 $\alpha,\beta,\gamma$ 是任意复数。级数 $F(\alpha,\beta,\gamma;z)$ 的收敛半径为 $1$，因此它在单位圆盘上定义一个全纯函数。这个全纯函数可以解析延拓为 $\C\setminus\{0,1\}$ 上的多值全纯函数。我们用同一个记号 $F(\alpha,\beta,\gamma;z)$ 表示这个多值全纯函数，并称它为超几何函数（hypergeometric function）。超几何函数 $F(\alpha,\beta,\gamma;z)$ 是超几何微分方程
\begin{equation}
z(1-z)\frac{d^2u}{dz^2}+(\gamma-(\alpha+\beta+1)z)\frac{du}{dz}-\alpha\beta u=0\tag{A.40}
\end{equation}
的解。超几何微分方程 (A.40) 有正则奇点 $0,1,\infty$。为简单起见，下面假定 $\gamma$ 不是负整数，也不是 $0$。在 $0$ 的邻域内，两个线性无关解由下式给出：
\begin{align}
u_{0,1}&=F(\alpha,\beta,\gamma;z),\nonumber\\
u_{0,2}&=z^{1-\gamma}F(1-\gamma+\alpha,1-\gamma+\beta,2-\gamma;z).\tag{A.41}
\end{align}
在 $1$ 的邻域内，两个线性无关解由下式给出：
\begin{align}
u_{1,1}&=F(\alpha,\beta,\alpha+\beta-\gamma+1;1-z)\tag{A.42}\\
u_{1,2}&=(1-z)^{\gamma-\alpha-\beta}F(\gamma-\alpha,\gamma-\beta,\gamma-\alpha-\beta;z).\tag{A.43}
\end{align}
类似地，在 $\infty$ 的邻域内，两个线性无关解由下式给出：
\begin{align}
u_{\infty,1}&=(-z)^{-\alpha}F\Bigl(\alpha,1+\alpha-\gamma,1+\alpha-\beta;\frac{1}{z}\Bigr)\tag{A.44}\\
u_{\infty,2}&=(-z)^{-\beta}F\Bigl(\beta,1+\beta-\gamma,1+\beta-\alpha;\frac{1}{z}\Bigr).\tag{A.45}
\end{align}

% [原文 p.162]
$F(\alpha,\beta,\gamma;z)$ 在点 $\infty$ 的邻域内的解析延拓由下式给出：
\begin{align}
F(\alpha,\beta,\gamma;z)&=\frac{\Gamma(\gamma)\Gamma(\beta-\alpha)}{\Gamma(\beta)\Gamma(\gamma-\alpha)}F\Bigl(\alpha,1+\alpha-\gamma,1+\alpha-\beta;\frac{1}{z}\Bigr)\nonumber\\
&\quad+\frac{\Gamma(\gamma)\Gamma(\alpha-\beta)}{\Gamma(\alpha)\Gamma(\gamma-\beta)}F\Bigl(\beta,1+\beta-\gamma,1+\beta-\alpha;\frac{1}{z}\Bigr).\tag{A.46}
\end{align}
其证明可在 Whittaker \& Watson [WW] 14.5 中找到。这一事实将用于计算定理 6.25 中的连接矩阵（connection matrix）。

% [Tutor Note: (a) $(\alpha)_n$ 即 Pochhammer 记号。原文说「对任意复数 $\alpha$」是对级数系数而言；一旦写成 $(\alpha)_n=\Gamma(\alpha+n)/\Gamma(\alpha)$，当 $\alpha$ 为非正整数时分子分母同时发散，这一等式应理解为极限（或用 $(\alpha)_n$ 的乘积定义直接取零）。为避免混淆，下文对 $\gamma$ 单独加了「不是负整数也不是 $0$」的假定。]
% [Tutor Note: (a) 收敛半径为 $1$？由相邻系数之比 $\frac{(\alpha+n)(\beta+n)}{(\gamma+n)(n+1)}\to1$ 即得。这同时说明一般情形下 $z=1$ 处级数发散，故 (A.39) 的收敛圆盘与 $1$ 处的解 (A.42) 并非同一展开。]
% [Tutor Note: (d) 这里的「多值」是实质性的：绕 $0$ 或绕 $1$ 解析延拓一周，$F$ 的值一般改变，改变量由 (A.41)–(A.46) 这类连接公式描述。定理 6.25 中的连接矩阵就是用 (A.46) 算出来的。]
% [Tutor Note: (b) (A.41) 中 $(0$ 处$)$ 的指数 $1-\gamma$ 来自指标方程：把 $u=z^{\rho}$ 代入 (A.40) 并取最低次项，得 $\rho(\rho-1)+\gamma\rho=0$，即 $\rho=0,1-\gamma$。$\gamma\notin\Z$ 正是两解线性无关的保证；$\gamma=1$ 时须改用含 $\log z$ 的解，原文用 $\gamma$ 的假定把这一情形排除。]
% [Tutor Note: (b) (A.42)(A.43) 中 $1$ 处的两个指数是 $0$ 与 $\gamma-\alpha-\beta$，同样由 (A.40) 在 $z=1$ 处的指标方程得到。注意 (A.42) 的第三个参数是 $\alpha+\beta-\gamma+1$（不是 $\gamma-\alpha-\beta$），它是把 $1-z$ 当作新变量后 (A.39) 的参数。]
% [Tutor Note: (b) (A.44)(A.45) 中 $\infty$ 处的两个指数是 $-\alpha,-\beta$。把 $z=1/t$ 代入 (A.40) 化成 $t=0$ 处的方程即可读出。因子 $(-z)^{-\alpha}$ 中的负号是为了固定分支：在 $z<0$ 的割缝上 $(-z)^{-\alpha}=|z|^{-\alpha}$。]
% [Tutor Note: (b) (A.46) 是连接公式。用 (A.44)(A.45) 可以自检系数：$z\to\infty$ 时右端第一项的行为是 $\frac{\Gamma(\gamma)\Gamma(\beta-\alpha)}{\Gamma(\beta)\Gamma(\gamma-\alpha)}(-z)^{-\alpha}$，第二项是 $\frac{\Gamma(\gamma)\Gamma(\alpha-\beta)}{\Gamma(\alpha)\Gamma(\gamma-\beta)}(-z)^{-\beta}$，与 (A.44)(A.45) 的首项一致；换言之 (A.46) 说的就是「$0$ 附近的解 $F(\alpha,\beta,\gamma;z)$ 在 $\infty$ 附近是 $u_{\infty,1},u_{\infty,2}$ 的线性组合」。]
% [Tutor Note: (a) (A.43) 中前因子 $(1-z)^{\gamma-\alpha-\beta}$ 的指数与第三个参数 $\gamma-\alpha-\beta$ 恰好相同，抄写时极易混为 $1$ 处的另一支；原文两处的确是同一个量。]
% [Tutor Note: (d) [WW] 指 Whittaker \& Watson 的《A Course of Modern Analysis》，其 14.5 节给出 (A.46) 的证明；本译本的参考文献中保留英文原条目，未翻书名。原文引用的「定理 6.25」在本书第 6 章，本块不涉及。]
% 索引（Index，原文 p.167）不译
